\documentclass{article}

\usepackage[preprint]{neurips_2026}
\usepackage[utf8]{inputenc}
\usepackage[T1]{fontenc}
\usepackage{hyperref}
\usepackage{url}
\usepackage{booktabs}
\usepackage{amsmath}
\usepackage{amsfonts}
\usepackage{nicefrac}
\usepackage{microtype}
\usepackage{xcolor}
\usepackage{graphicx}
\usepackage{float}
\hypersetup{hidelinks}

\title{Segmentation-Aware Machine Learning for Intercity Travel Mode Choice Prediction}

\author{%
  CS439 Spring 2026\\
  Student\\
  \texttt{ud38@scarletmail.rutgers.edu}\\
}

\begin{document}
\maketitle

\begin{abstract}
Reliable travel mode prediction helps transportation planners understand how travelers trade off cost, time, party size, and household income when choosing among air, train, bus, and car.  This project studies the Travel Mode Choice dataset, which contains 840 alternative-level observations from 210 intercity travelers, and frames prediction as a leakage-sensitive grouped choice problem: each traveler contributes four alternatives, exactly one of which is selected.  I propose a lightweight hybrid pipeline that first clusters travelers using K-Means on choice-set attributes and then adds the learned segment labels to supervised classifiers.  The pipeline includes strict group-level train/test separation, binary mapping, one-hot encoding, standardized numerical features, utility-inspired feature engineering, and multiple baselines.  On a held-out set of 53 unseen travelers, the best operational model, a segmentation-aware Random Forest, achieves 94.3\% top-1 choice-set accuracy and 96.2\% top-2 accuracy, substantially improving over logistic regression (69.8\% top-1) and rule-based baselines (26.4--32.1\% top-1).  XGBoost obtains the strongest row-level discrimination with ROC-AUC 0.985 and F1 0.925.  TreeSHAP analysis indicates that relative travel time, terminal waiting time, in-vehicle time, and generalized cost are the dominant predictors, aligning the learned models with transportation intuition.
\end{abstract}

\paragraph{Code repository.} \url{https://github.com/utkarshdubey/cs439-mode-choice-final}.  The accompanying reproducibility package contains the full preprocessing, segmentation, model training, evaluation, and figure-generation pipeline.

\section{Introduction}
Mode choice modeling asks which transportation alternative a traveler will select from a feasible set.  Although classical transportation planning often relies on discrete-choice models, modern data science workflows can complement these models by learning nonlinear interactions, producing operational predictions, and exposing feature-level explanations.  In this project, the application is intercity travel between Australian cities, where each traveler chooses one of four modes: air, train, bus, or car.  The practical objective is not only to estimate whether a single row is chosen, but also to rank the four alternatives for each unseen traveler and select the most likely mode.

The problem is deceptively easy to evaluate incorrectly.  The dataset is stored in long format, so each traveler appears in four rows.  A standard random row split would place some alternatives from the same traveler in training and others in testing, leaking traveler-specific income, party size, and choice-set structure into evaluation.  Therefore, this project treats the traveler as the atomic unit for splitting, clustering, cross-validation, and final testing.

The main research question is: \emph{Can traveler segmentation and utility-inspired feature engineering improve interpretable machine learning for intercity travel mode choice under a leakage-safe evaluation protocol?}  The project makes three contributions.  First, it builds a rigorous preprocessing pipeline that separates train and test travelers, maps binary outcomes, one-hot encodes categorical variables, scales numerical variables, and checks data integrity.  Second, it combines unsupervised K-Means traveler segmentation with supervised prediction, allowing the final models to condition on both alternative-level attributes and traveler cohort structure.  Third, it evaluates logistic regression, Random Forest, Extra Trees, and XGBoost against rule-based baselines using both row-level metrics and choice-set metrics, and it interprets the strongest model using TreeSHAP.

\section{Related Work}
\textbf{Discrete-choice modeling.}  Travel mode choice is a canonical application of random utility theory and multinomial or conditional logit models.  McFadden's conditional logit formulation provides a probabilistic model of discrete alternatives by representing utility as observed covariates plus random error \citep{mcfadden1974}.  Train's treatment of discrete choice methods further formalizes estimation, simulation, and interpretation for transportation and economics applications \citep{train2009}.  These methods remain valuable because their coefficients have clear behavioral meaning, but the linear-in-parameters utility assumption can limit predictive performance when interactions are complex.

\textbf{Machine learning for mode choice.}  Recent studies compare traditional logit models with machine learning classifiers for travel behavior prediction.  Hagenauer and Helbich evaluate several classifiers for travel mode choice and show the value of comparing predictive accuracy with interpretable variable importance \citep{hagenauer2017}.  Ali et al. similarly emphasize that machine learning can improve prediction while discrete-choice models remain useful for interpretation \citep{ali2021}.  This project follows that hybrid spirit: logistic regression is retained as a transparent baseline, while tree ensembles are used to capture nonlinear interactions.

\textbf{Segmentation and explainability.}  K-Means clustering is a standard unsupervised method for partitioning observations into compact groups by minimizing within-cluster squared distances \citep{lloyd1982}.  Tree ensembles such as Random Forests \citep{breiman2001} and gradient-boosted trees \citep{chen2016xgboost} frequently perform well on tabular data.  However, higher predictive power can reduce transparency, so post-hoc explanation methods are useful.  SHAP provides additive feature attributions grounded in Shapley values and is widely used to explain tree-based models \citep{lundberg2017}.  This project uses K-Means for cohort discovery and TreeSHAP to audit whether the final models rely on plausible transportation features.

\section{Methodology}
\subsection{Dataset and problem formulation}
The dataset is loaded from \texttt{statsmodels.datasets.modechoice}, whose documentation describes it as a 1987 intercity mode-choice sample with 210 travelers and four alternatives per traveler \citep{statsmodels_modechoice}.  The raw fields are traveler ID, mode ID, binary choice label, terminal waiting time (\texttt{ttme}), in-vehicle cost (\texttt{invc}), in-vehicle time (\texttt{invt}), generalized cost (\texttt{gc}), household income (\texttt{hinc}), and party size (\texttt{psize}).  I map mode IDs to \{air, train, bus, car\} and cast the target \texttt{choice} into a binary variable $y_{ij}\in\{0,1\}$ for traveler $i$ and mode $j$.

The row-level model estimates $\hat{p}_{ij}=P(y_{ij}=1\mid x_{ij})$.  For operational choice prediction, probabilities are converted into a traveler-level decision by selecting
\begin{equation}
\hat{j}_i = \arg\max_{j\in\{\text{air, train, bus, car}\}} \hat{p}_{ij}.
\end{equation}
The primary metric is top-1 choice-set accuracy, $\frac{1}{N}\sum_i \mathbb{1}[\hat{j}_i=j_i]$, because a transportation planner ultimately needs the chosen mode within each traveler choice set.  I also report top-2 accuracy, row-level F1, ROC-AUC, PR-AUC, and log loss.

\subsection{Preprocessing and leakage prevention}
The preprocessing pipeline is deliberately conservative.  I first verify that the dataset has no missing values, no duplicate rows, exactly four alternatives per traveler, and exactly one chosen alternative per traveler.  The final split uses \texttt{GroupShuffleSplit} with traveler ID as the group, producing 157 training travelers (628 rows) and 53 held-out test travelers (212 rows).  Cross-validation uses five-fold \texttt{GroupKFold}, again grouped by traveler ID.

Feature transformations are fit only on training data inside scikit-learn pipelines \citep{pedregosa2011}.  Binary mapping is used for the target and for two engineered indicators: \texttt{is\_lowest\_gc} and \texttt{is\_fastest}.  One-hot encoding is used for \texttt{mode\_name} and the K-Means segment label.  Standard scaling is applied to numerical features before logistic regression, clustering, and tree pipelines for a uniform preprocessing interface.  The engineered variables are utility-inspired: total time $=\texttt{ttme}+\texttt{invt}$, cost per income, party-adjusted cost, difference from the minimum generalized cost in the choice set, difference from the fastest alternative, and ratios to choice-set mean cost and mean time.  These features use only covariates available within the traveler's alternatives and never use the label.

\subsection{Traveler segmentation}
To add interpretable cohort information, I build a traveler-level matrix using income, party size, and summary statistics for each traveler's four alternatives.  For each of \texttt{gc}, \texttt{ttme}, \texttt{invc}, \texttt{invt}, and total time, the matrix includes minimum, mean, maximum, range, and mode-specific values.  K-Means is fit on training travelers only, and held-out travelers are assigned to the nearest training-fitted centroid.  I test $k\in\{2,3,4,5,6\}$ and select $k=2$ because it yields the highest silhouette score and lowest Davies-Bouldin score in the training data.  The two segments correspond to longer, higher generalized-cost choice sets versus shorter, lower-cost choice sets.

\subsection{Supervised models and baselines}
The supervised baselines include logistic regression with raw features, logistic regression with engineered features, and logistic regression with engineered features plus segment labels.  The tree-based models are Random Forest, Extra Trees, and XGBoost.  Class weights or positive-class scaling are used because each four-row choice set has one positive and three negative rows.  Hyperparameters are kept intentionally small and stable to avoid overfitting the 210-traveler sample: Random Forest and Extra Trees use 300 trees, maximum depth 5, and minimum leaf size 3; XGBoost uses 100 trees of depth 2 with learning rate 0.08.  Three non-learning baselines are also included: always choose the training-popular mode, choose the alternative with lowest generalized cost, and choose the fastest total-time alternative.

\section{Experiments \& Results}
\subsection{Segmentation diagnostics}
Table~\ref{tab:segments} summarizes the selected clusters.  Segment 0 contains 110 travelers and has substantially higher mean generalized cost and total time.  Segment 1 contains 100 travelers and represents shorter, lower-cost choice sets.  Figure~\ref{fig:segments} shows that the two clusters separate clearly in the first two PCA dimensions, supporting their use as interpretable cohort labels.

\begin{table}[H]
\centering
\small
\caption{K-Means segment summaries computed after assigning all travelers with centroids fit on training travelers only.}
\label{tab:segments}
\begin{tabular}{lrrrrr}
\toprule
Segment & Travelers & Mean income & Mean party size & Mean generalized cost & Mean total time \\
\midrule
0 & 110 & 35.34 & 1.73 & 146.70 & 688.31 \\
1 & 100 & 33.68 & 1.76 & 71.47 & 336.45 \\
\bottomrule
\end{tabular}
\end{table}

\begin{figure}[H]
\centering
\includegraphics[width=0.62\linewidth]{./pca_segments.pdf}
\caption{PCA projection of traveler-level K-Means segments.  Segment 0 represents longer and more expensive choice sets, while Segment 1 represents shorter and lower-cost choice sets.}
\label{fig:segments}
\end{figure}

\subsection{Predictive performance}
Table~\ref{tab:results} reports the main supervised results on the held-out travelers.  The rule-based baselines are weak: lowest generalized cost obtains 32.1\% top-1 choice-set accuracy, and fastest total time obtains 26.4\%.  Logistic regression with raw features reaches 58.5\% top-1 accuracy.  Adding utility-inspired engineered features raises logistic top-1 accuracy to 69.8\%, showing that the feature design contributes more than the model family alone.  The best operational model is Random Forest with segment labels, which reaches \textbf{94.3\%} top-1 accuracy and 96.2\% top-2 accuracy.  XGBoost is the strongest row-level ranker, with \textbf{0.985 ROC-AUC}, \textbf{0.966 PR-AUC}, and \textbf{0.925 F1}, but its top-1 choice-set accuracy is slightly below Random Forest.

\begin{table}[H]
\centering
\small
\caption{Held-out performance on 53 unseen travelers.  CV ROC-AUC is from five-fold GroupKFold on training travelers.}
\label{tab:results}
\resizebox{\linewidth}{!}{%
\begin{tabular}{lrrrrrr}
\toprule
Model & CV ROC-AUC & F1 & ROC-AUC & PR-AUC & Top-1 choice acc. & Top-2 acc. \\
\midrule
Logit-basic & 0.797 $\pm$ 0.049 & 0.500 & 0.707 & 0.641 & 0.585 & 0.736 \\
Logit-engineered & 0.885 $\pm$ 0.034 & 0.644 & 0.838 & 0.762 & 0.698 & 0.849 \\
Logit+segments & 0.886 $\pm$ 0.036 & 0.639 & 0.835 & 0.762 & 0.698 & 0.849 \\
RandomForest+segments & 0.966 $\pm$ 0.020 & 0.923 & 0.972 & 0.952 & \textbf{0.943} & \textbf{0.962} \\
ExtraTrees+segments & 0.926 $\pm$ 0.020 & 0.723 & 0.911 & 0.858 & 0.755 & \textbf{0.962} \\
XGBoost+segments & \textbf{0.972 $\pm$ 0.023} & \textbf{0.925} & \textbf{0.985} & \textbf{0.966} & 0.925 & \textbf{0.962} \\
\bottomrule
\end{tabular}}
\end{table}

\begin{figure}[H]
\centering
\includegraphics[width=0.94\linewidth]{./choice_set_accuracy.pdf}
\caption{Top-1 choice-set accuracy for rules and supervised models.  The ensemble models substantially outperform utility-only rules and logistic regression.}
\label{fig:accuracy}
\end{figure}

\subsection{Error analysis and interpretation}
The Random Forest confusion matrix in Figure~\ref{fig:cm} shows that the best operational model correctly predicts most modes and makes only three traveler-level mistakes on the 53-person test set.  The remaining errors are plausible confusions between alternatives with similar time-cost profiles.  ROC curves in Figure~\ref{fig:roc} confirm that the tree ensembles dominate logistic regression as row-level rankers, while XGBoost has the best discrimination across thresholds.

\begin{figure}[H]
\centering
\includegraphics[width=0.48\linewidth]{./rf_confusion_matrix.pdf}
\caption{Choice-set confusion matrix for RandomForest+segments. Rows are true chosen modes and columns are predicted modes.}
\label{fig:cm}
\end{figure}

\begin{figure}[H]
\centering
\includegraphics[width=0.68\linewidth]{./roc_curves.pdf}
\caption{ROC curves on held-out travelers.  XGBoost gives the highest row-level ROC-AUC, while Random Forest gives the best top-1 choice-set accuracy.}
\label{fig:roc}
\end{figure}

TreeSHAP results in Figure~\ref{fig:shap} provide a sanity check on model behavior.  The top features are relative total time within the choice set, terminal waiting time, in-vehicle time, generalized cost, car and train indicators, in-vehicle cost, and generalized-cost gap.  This supports the claim that the high-performing model is not exploiting arbitrary identifiers; it relies on transportation-relevant variables that match the expected cost-time tradeoff.

\begin{figure}[H]
\centering
\includegraphics[width=0.62\linewidth]{./shap_top_features.pdf}
\caption{Top TreeSHAP features for XGBoost.  Relative time and cost variables dominate, which is behaviorally plausible for intercity mode choice.}
\label{fig:shap}
\end{figure}

\section{Discussion and Limitations}
The results show that careful evaluation design matters as much as model choice.  If alternatives from the same traveler were randomly split across train and test sets, evaluation would be inflated by leakage.  Grouped splitting gives a harder and more realistic test: the model must generalize to entirely unseen travelers.  Under that protocol, feature engineering produces a clear improvement over raw-feature logistic regression, and tree ensembles capture interactions that logistic regression misses.  The segmentation step is most useful for interpretation and cohort auditing; it does not dramatically change logistic regression accuracy, but it gives a concise way to distinguish long/high-cost and short/low-cost traveler contexts.

The main limitation is dataset size.  The sample contains only 210 travelers and is choice-based, with oversampling of less popular modes, so the model should not be interpreted as estimating real population mode shares without additional weighting or external validation.  The dataset also lacks variables that are important in modern travel behavior, such as departure time, service reliability, comfort, vehicle ownership, emissions preferences, and geographic context.  Finally, the independent binary formulation ranks alternatives well, but a true conditional logit or neural choice model could explicitly normalize probabilities within each choice set.

\section{Conclusion}
This project developed a reproducible, leakage-safe data science pipeline for intercity travel mode choice prediction.  The strongest operational model, a segmentation-aware Random Forest, achieved 94.3\% top-1 choice-set accuracy on unseen travelers, while XGBoost achieved the best row-level ROC-AUC of 0.985.  The analysis demonstrates that a low-effort but rigorous hybrid pipeline--grouped splitting, utility-inspired features, K-Means traveler segmentation, tree ensembles, and SHAP interpretation--can satisfy both predictive and interpretability goals.  Future work should validate the pipeline on larger, more recent travel surveys, incorporate sample weights for population inference, and compare the binary ranking setup with dedicated multinomial and conditional logit architectures.

\begin{thebibliography}{10}

\bibitem[Ali et~al.(2021)]{ali2021}
Ali, N.~F.~M., Sadullah, A.~F.~M., Majeed, A.~P.~P.~A., Razman, M.~A.~M., Zakaria, M.~F., and Nasir, M.~A.
\newblock Travel mode choice modeling: Predictive efficacy between machine learning models and discrete choice model.
\newblock \emph{The Open Transportation Journal}, 15:241--257, 2021.

\bibitem[Breiman(2001)]{breiman2001}
Breiman, L.
\newblock Random forests.
\newblock \emph{Machine Learning}, 45:5--32, 2001.

\bibitem[Chen and Guestrin(2016)]{chen2016xgboost}
Chen, T. and Guestrin, C.
\newblock XGBoost: A scalable tree boosting system.
\newblock In \emph{Proceedings of the 22nd ACM SIGKDD International Conference on Knowledge Discovery and Data Mining}, pages 785--794, 2016.

\bibitem[Hagenauer and Helbich(2017)]{hagenauer2017}
Hagenauer, J. and Helbich, M.
\newblock A comparative study of machine learning classifiers for modeling travel mode choice.
\newblock \emph{Expert Systems with Applications}, 78:273--282, 2017.

\bibitem[Lloyd(1982)]{lloyd1982}
Lloyd, S.
\newblock Least squares quantization in PCM.
\newblock \emph{IEEE Transactions on Information Theory}, 28(2):129--137, 1982.

\bibitem[Lundberg and Lee(2017)]{lundberg2017}
Lundberg, S.~M. and Lee, S.-I.
\newblock A unified approach to interpreting model predictions.
\newblock In \emph{Advances in Neural Information Processing Systems}, volume 30, 2017.

\bibitem[McFadden(1974)]{mcfadden1974}
McFadden, D.
\newblock Conditional logit analysis of qualitative choice behavior.
\newblock In Zarembka, P., editor, \emph{Frontiers in Econometrics}, pages 105--142. Academic Press, 1974.

\bibitem[Pedregosa et~al.(2011)]{pedregosa2011}
Pedregosa, F., Varoquaux, G., Gramfort, A., Michel, V., Thirion, B., Grisel, O., Blondel, M., Prettenhofer, P., Weiss, R., Dubourg, V., VanderPlas, J., Passos, A., Cournapeau, D., Brucher, M., Perrot, M., and Duchesnay, E.
\newblock Scikit-learn: Machine learning in Python.
\newblock \emph{Journal of Machine Learning Research}, 12:2825--2830, 2011.

\bibitem[Statsmodels Developers(2025)]{statsmodels_modechoice}
Statsmodels Developers.
\newblock Travel Mode Choice dataset documentation.
\newblock \url{https://www.statsmodels.org/stable/datasets/generated/modechoice.html}, accessed 2026.

\bibitem[Train(2009)]{train2009}
Train, K.~E.
\newblock \emph{Discrete Choice Methods with Simulation}.
\newblock Cambridge University Press, 2nd edition, 2009.

\end{thebibliography}

\end{document}
